\subsection{INTEGRALS}

We have obtained (II, p.~54, th. 2) a primitive of a regulated function on an interval I as the uniform limit of primitives of step functions. This procedure can be expressed in a slightly different way: let \(x_0\) , \(x\) be two arbitrary points of I such that \(x_0 < x\) ; we call a subdivision of the interval \([x_0, x]\) any sequence of intervals \([x_i, x_{i+1}]\) with union \([x_0, x]\) , where \((x_i)_{0 \leqslant i \leqslant n}\) is a strictly increasing sequence of points of \([x_0, x]\) such that \(x_n = x\) . We shall call a Riemann sum, relative to a vector function \(\mathbf{f}\) defined on I and to the subdivision formed by the \([x_i, x_{i+1}]\) , any expression of the form \(\sum_{i=0}^{n-1} \mathbf{f}(t_i)(x_{i+1} - x_i)\) where the \(t_i\) belong to \([x_i, x_{i+1}]\) for \(0 \leqslant i \leqslant n-1\) . One then has the following proposition:

PROPOSITION 5. Let \(\mathbf{f}\) be a regulated function on an interval I, let \(\mathbf{g}\) be a primitive of \(\mathbf{f}\) on I, and \([x_0, x]\) a compact interval contained in I. For every \(\varepsilon > 0\) there exists a number \(\rho > 0\) such that for every subdivision of \([x_0, x]\) into intervals of length \(\leqslant \rho\) one has

\[
\left\| \mathbf {g} (x) - \mathbf {g} (x _ {0}) - \sum_ {i = 0} ^ {n - 1} \mathbf {f} (t _ {i}) (x _ {i + 1} - x _ {i}) \right\| \leqslant \varepsilon\tag{1}
\]

for every Riemann sum relative to this subdivision.

Indeed, let \(\mathbf{f}_1\) be a step function such that \(\| \mathbf{f}(y) - \mathbf{f}_1(y)\| \leqslant \varepsilon\) for every \(y\in [x_0,x]\) ; one has, denoting a primitive of \(\mathbf{f}_1\) on I by \(\mathbf{g}_1\) ,

\[
\left\| \mathbf {g} (x) - \mathbf {g} \left(x _ {0}\right) - \left(\mathbf {g} _ {1} (x) - \mathbf {g} _ {1} \left(x _ {0}\right)\right) \right\| \leqslant \varepsilon \left(x - x _ {0}\right)
\]

by the mean value theorem, and on the other hand

\[
\left\| \sum_ {i = 0} ^ {n - 1} \mathbf {f} (t _ {i}) (x _ {i + 1} - x _ {i}) - \sum_ {i = 0} ^ {n - 1} \mathbf {f} _ {1} (t _ {i}) (x _ {i + 1} - x _ {i}) \right\| \leqslant \varepsilon (x - x _ {0}).
\]

It thus suffices to prove the proposition when \(\mathbf{f}\) is a step function. Let \((y_{k})_{1\leqslant k\leqslant m}\) be the strictly increasing finite sequence of points of discontinuity of \(\mathbf{f}\) in \([x_0, x]\) . For every subdivision of \([x_0, x]\) into intervals of length \(\leqslant \rho\) each of the points \(y_{k}\) belongs to at most two of these intervals; there can therefore be no more than \(2m\) intervals on which \(\mathbf{f}\) is not constant; but, on such an interval \([x_i, x_{i+1}]\) one has

\[
\left\| \mathbf {g} \left(x _ {i + 1}\right) - \mathbf {g} \left(x _ {i}\right) - \mathbf {f} \left(t _ {i}\right) \left(x _ {i + 1} - x _ {i}\right) \right\| \leqslant 2 \mathrm{M} \left(x _ {i + 1} - x _ {i}\right)
\]

on denoting by M the least upper bound of \(\|f\|\) on \([x_{0}, x]\) ; on the other hand, when f is constant on \([x_{i}, x_{i+1}]\) one has

\[
\mathbf {g} (x _ {i + 1}) - \mathbf {g} (x _ {i}) - \mathbf {f} (t _ {i}) (x _ {i + 1} - x _ {i}) = 0.
\]

One thus sees that the difference \(\left\|\mathbf{g}(x)-\mathbf{g}(x_{0})-\sum_{i=0}^{n-1}\mathbf{f}(t_{i})(x_{i+1}-x_{i})\right\|\) cannot exceed \(4Mm\rho\) ; it therefore suffices to take \(\rho\leqslant\varepsilon/4Mm\) to obtain (1).

Remark. 1) When \(\mathbf{f}\) is continuous prop. 5 can be proved more simply: since \(\mathbf{f}\) is uniformly continuous on \([x_0, x]\) there exists a \(\rho > 0\) such that on every interval of length \(\leqslant \rho\) contained in \([x_0, x]\) the oscillation of \(\mathbf{f}\) is \(\leqslant \frac{\varepsilon}{x - x_0}\) ; for every subdivision of \([x_0, x]\) into intervals \([x_i, x_{i+1}]\) of length \(\leqslant \rho\) and every choice of \(t_i\) in \([x_i, x_{i+1}]\) for \(0 \leqslant i \leqslant n-1\) the step function \(\mathbf{f}_1\) equal to \(\mathbf{f}(t_i)\) on \([x_i, x_{i+1}]\) ( \(0 \leqslant i \leqslant n-1\) ), and to \(\mathbf{f}(x)\) at the point \(x\) , is such that \(\| \mathbf{f}(y) - \mathbf{f}_1(y)\| \leqslant \frac{\varepsilon}{x - x_0}\) on \([x_0, x]\) ; if \(\mathbf{g}_1\) is a primitive of \(\mathbf{f}_1\) one has \(\mathbf{g}_1(x) - \mathbf{g}_1(x_0) = \sum_{i=0}^{n-1} \mathbf{f}(t_i)(x_{i+1} - x_i)\) , so the relation (1) follows immediately from the mean value theorem.

In the rest of this chapter we shall confine ourselves to the study of primitives of regulated functions on an interval I. For such a function \(\mathbf{f}\) , with values in E, a primitive \(\mathbf{g}\) of \(\mathbf{f}\) , and for two arbitrary points \(x_0\) , \(x\) of I, the element \(\mathbf{g}(x) - \mathbf{g}(x_0)\) of E (which is clearly the same, no matter which primitive \(\mathbf{g}\) of \(\mathbf{f}\) one considers) is called the integral of the function \(\mathbf{f}\) from \(x_0\) to \(x\) (or over the compact interval \([x_0, x]\) ) and is denoted by \(\int_{x_0}^{x} \mathbf{f}(t) dt\) or \(\int_{x_0}^{x} \mathbf{f}\) . This name and notation have their origin in prop. 5 of II, p.~57, which shows that an integral can be approximated arbitrarily closely by a Riemann sum; more particularly, one can write, taking subdivisions of \([x_0, x]\) into equal intervals,

\[
\frac {1}{x - x _ {0}} \int_ {x _ {0}} ^ {x} \mathbf {f} (t) d t = \lim _ {n \rightarrow \infty} \frac {1}{n} \sum_ {k = 0} ^ {n - 1} \mathbf {f} \left(x _ {0} + k \frac {x - x _ {0}}{n}\right).\tag{2}
\]

In other words, the element \(\frac{1}{x-x_{0}}\int_{x_{0}}^{x}\mathbf{f}(t)dt\) is the limit of the arithmetic means of the values of f at the left-hand endpoints of the intervals of a subdivision of \([x_{0},x]\) into equal intervals; one also calls it the mean (or mean value) of the function f on the interval \([x_{0},x]\) .

By definition, the function \(x \mapsto \int_{x_{0}}^{x} \mathbf{f}(t) dt\) is none other than the primitive of f which vanishes at the point \(x_{0} \in I\) ; one also denotes it by \(\int_{x_{0}} \mathbf{f}(t) dt\) or \(\int_{x_{0}} f\) .

Remarks. 2) For an arbitrary function \(\mathbf{h}\) defined on I, with values in E, the element \(\mathbf{h}(x) - \mathbf{h}(x_0)\) is also written as \(\left.\mathbf{h}(t)\right|_{x_0}^x\) ; with this notation one sees that if \(\mathbf{g}\) is any primitive of a regulated function \(\mathbf{f}\) on I, one has

\[
\int_ {x _ {0}} ^ {x} \mathbf {f} (t) d t = \mathbf {g} (t) \left| _ {x _ {0}} ^ {x}. \right.\tag{3}
\]

\begin{enumerate}
\def\labelenumi{\arabic{enumi})}
\setcounter{enumi}{2}
\item
  The expressions \(\int_{x_{0}}^{x}\mathbf{f}(t)dt\) and \(\mathbf{g}(t)\bigg|_{x_{0}}^{x}\) are abbreviating symbols representing assemblies in which the letters x, \(x_{0}\) , f, g, but not the letter t, appear (cf.~Set Theory, I, p.~15); one says that among these symbols t is a ``dummy variable''; one can thus replace t by any other variable distinct from x, \(x_{0}\) , f and g (and from variables which may possibly enter into the proof where these symbols appear) without changing the sense of the symbol so obtained (the reader may compare these symbols with symbols such as \(\sum_{i=1}^{n}x_{i}\) , or \(\bigcup_{i}X_{i}\) , where i is likewise a dummy variable).
\item
  The approximation of an integral by Riemann sums is closely connected to one of the historical origins of the concept of an integral, the problem of the measure of areas. We shall come back to this point in the Book on Integration which is devoted to generalizations of the concept of integral to which this problem leads; in these generalizations the functions to be ``integrated'' are not necessarily defined on a subset of R; moreover, even when one deals with (not necessarily regulated) real functions f of a real variable for which one can define an integral \(\int_{x_{0}}^{x}f(t)dt\) , the function \(x\mapsto\int_{x_{0}}^{x}f(t)dt\) is not always a primitive of f, and there exist functions which have a primitive but are not ``integrable'' in the sense to which we allude.
\end{enumerate}
% semantic-fragments: legacy-input-seam
\relax{}
